% ============================================================================
%  Capitolo 24 -- Oltre un solo core                             (Lezione 24)
% ============================================================================
\chapter{Oltre un solo core: RISC-V superscalare, HPC, GPU e FPGA}
\label{ch:beyond}

\epigraphline{Se doveste arare un campo, che cosa preferireste usare: due buoi
robusti o 1024 polli?}{attribuita a Seymour Cray}

\begin{chapgoals}
Alla fine di questo capitolo saprete
\begin{itemize}
\item scrivere l'equazione delle prestazioni e dire quale fattore attacca ciascuna tecnica;
\item spiegare come un processore superscalare fuori ordine esegue più
  istruzioni per ciclo, e nominare core \riscv che lo fanno;
\item leggere un ciclo scritto con l'estensione vettoriale di \riscv;
\item descrivere un supercalcolatore, calcolare uno speed-up con la legge di Amdahl, e dire perché il limite è l'energia;
\item dire in che cosa sono bravi una GPU e una FPGA, e che cosa costa ciascuna al programmatore.
\end{itemize}
\end{chapgoals}

Questo capitolo conclusivo è uno sguardo in avanti. Parte da un ciclo di
\vermone e lo segue su quattro tipi di macchina: le stesse istruzioni, su
macchine che ne eseguono molte insieme. Non fa parte della prova intercorso. I
listati vettoriali \riscv sono stati provati su QEMU; i frammenti CUDA e
Verilog sono forme standard da libro di testo che per questo libro non sono
state compilate. Le cifre sui supercalcolatori sono arrotondate e vengono dalle
liste TOP500 dal 2022 al 2025.

% ----------------------------------------------------------------------------
\section{Riscaldamento sulla macchina}

\asmfile[caption={Il ciclo che svuota la mappa, dalla \code{new\_game} di \vermone.},label=lst:by:loop]{code/l24/c03.S}

Il Listato~\ref{lst:by:loop} pone a zero 960 byte: quattro istruzioni per byte,
3\,840 istruzioni in tutto, un byte alla volta e un'istruzione alla volta. Nel
caso migliore un'istruzione per ciclo di clock e 160 milioni di cicli al
secondo sul kit: circa 24 microsecondi. A nessuno serve che questo ciclo sia
più veloce; ma è abbastanza piccolo da poter essere seguito su ogni macchina di
questo capitolo. Quanti modi ci sono per renderlo più veloce?

\subsection{L'equazione delle prestazioni}

\begin{theorem}{Equazione delle prestazioni}
Il tempo che un programma impiega è il prodotto di tre numeri,
\[ T = N\times \mathit{CPI}\times t , \]
dove $N$ è il numero di istruzioni eseguite, $\mathit{CPI}$ il numero medio di
cicli di clock per istruzione e $t$ la durata di un ciclo di clock.
\end{theorem}

Per il ciclo precedente sul kit: $N=3\,840$, circa 1 ciclo ciascuna, 6,25
nanosecondi per ciclo. Ogni fattore può essere attaccato: \emph{meno
istruzioni}, con algoritmi migliori, istruzioni più ricche, vettori;
\emph{meno cicli per istruzione}, con le pipeline e l'esecuzione superscalare;
\emph{un ciclo più breve}, con transistor più veloci e pipeline più profonde.
E quando un processore non basta: più processori che lavorano insieme. Ogni
macchina di questo capitolo è una risposta diversa alla stessa equazione
\citep{patterson2020}.

\subsection{Perché non semplicemente un clock più veloce?}

Per trent'anni ogni nuova generazione di transistor è stata più piccola, più
veloce, e non più calda per millimetro quadrato. Intorno al 2005 questo è
finito: transistor più piccoli arrivano ancora, ma un clock più veloce
significa più calore di quanto un chip possa smaltire
(Capitolo~\ref{ch:electronics}), e da allora le frequenze di clock sono rimaste
a pochi gigahertz. La legge di Moore, più transistor, è andata avanti; la
regola che teneva costante la densità di potenza, nota come scaling di Dennard,
è finita. I transistor in più sono andati altrove: processori più larghi, più
core, unità vettoriali, acceleratori. Le prestazioni vengono ora dal fare molte
cose insieme: dal \term{parallelismo}. I 160~MHz del kit sono lenti per scelta:
costano milliwatt.

% ----------------------------------------------------------------------------
\section{Processori superscalari}

\subsection{Da una alla volta a molte insieme}

\begin{table}[htb]
\centering\small
\caption{Da un'istruzione alla volta a molte insieme.}\label{tab:by:kinds}
\begin{tabularx}{\textwidth}{@{}P{0.21\textwidth}LP{0.14\textwidth}L@{}}
\toprule
\textbf{Processore} & \textbf{Come funziona} & \textbf{Istr.\ per ciclo} & \textbf{Esempio}\\
\midrule
Multiciclo & un'istruzione si completa prima che cominci la successiva & meno di 1 & il modello del Capitolo~\ref{ch:processor}\\
Con pipeline & più istruzioni in volo, ciascuna in uno stadio diverso & al più 1 & l'ESP32-C6 del kit: quattro stadi\\
Superscalare, in ordine & due o più pipeline fianco a fianco & fino a 2 o più & piccoli core applicativi\\
Superscalare, fuori ordine & le istruzioni partono appena i loro operandi sono pronti & fino a 4, 6, 8, \ldots & portatili, telefoni, server\\
\bottomrule
\end{tabularx}
\end{table}

Il Capitolo~\ref{ch:instructions} ha introdotto la pipeline e nominato i
processori superscalari in una riga; questo paragrafo apre quella riga
(Tabella~\ref{tab:by:kinds}). La \emph{larghezza} di un processore è il numero
di istruzioni che può avviare in un ciclo. Con due pipeline fianco a fianco e i
cinque stadi classici, sei istruzioni si completano in sette cicli
(Tabella~\ref{tab:by:two}); a una sola pipeline ne servono dieci, e a un
processore senza pipeline trenta.

\begin{table}[htb]
\centering\small
\caption{Due pipeline fianco a fianco: sei istruzioni in sette cicli, se niente
si mette di mezzo.}\label{tab:by:two}
\begin{tabular}{@{}l*{7}{>{\centering\arraybackslash}p{8mm}}@{}}
\toprule
\textbf{ciclo} & 1 & 2 & 3 & 4 & 5 & 6 & 7\\
\midrule
$i_1$ & \cellcolor{prgaccent!20}IF & \cellcolor{prgamber!25}ID & \cellcolor{prgviolet!18}EX & \cellcolor{prggreen!20}MEM & \cellcolor{prgred!15}WB &&\\
$i_2$ & \cellcolor{prgaccent!20}IF & \cellcolor{prgamber!25}ID & \cellcolor{prgviolet!18}EX & \cellcolor{prggreen!20}MEM & \cellcolor{prgred!15}WB &&\\
$i_3$ && \cellcolor{prgaccent!20}IF & \cellcolor{prgamber!25}ID & \cellcolor{prgviolet!18}EX & \cellcolor{prggreen!20}MEM & \cellcolor{prgred!15}WB &\\
$i_4$ && \cellcolor{prgaccent!20}IF & \cellcolor{prgamber!25}ID & \cellcolor{prgviolet!18}EX & \cellcolor{prggreen!20}MEM & \cellcolor{prgred!15}WB &\\
$i_5$ &&& \cellcolor{prgaccent!20}IF & \cellcolor{prgamber!25}ID & \cellcolor{prgviolet!18}EX & \cellcolor{prggreen!20}MEM & \cellcolor{prgred!15}WB\\
$i_6$ &&& \cellcolor{prgaccent!20}IF & \cellcolor{prgamber!25}ID & \cellcolor{prgviolet!18}EX & \cellcolor{prggreen!20}MEM & \cellcolor{prgred!15}WB\\
\bottomrule
\end{tabular}
\end{table}

\subsection{Che cosa si mette di mezzo: le dipendenze}

\asmfile[caption={Dipendenze tra istruzioni vicine.},label=lst:by:dep]{code/l24/c04.S}

Questo è il caso ideale, ed è raro. Nel Listato~\ref{lst:by:dep} l'istruzione
$i_2$ non può essere eseguita accanto a $i_1$, perché ne usa il risultato: una
\term{dipendenza sui dati}. $i_3$ non dipende da nessuno: perché dovrebbe
aspettare? E dopo il salto $i_4$ il processore non sa che cosa viene dopo: una
\term{dipendenza di controllo}. Il codice reale è pieno di catene di questo
tipo, e la sola larghezza non basta. Un processore superscalare che lavora in
ordine si blocca a $i_2$, e tutto ciò che sta dietro aspetta. Il rimedio è
lasciare che $i_3$ sorpassi.

\subsection{L'esecuzione fuori ordine}

\begin{enumerate}
\item \textbf{Prelievo e decodifica, in ordine.} Più istruzioni per ciclo
  entrano in una finestra di decine o centinaia di istruzioni.
\item \textbf{Ridenominazione.} Ogni risultato riceve un registro fisico nuovo,
  così che due usi del nome \reg{t0} non collidono più.
\item \textbf{Esecuzione quando è pronta.} Un'istruzione parte appena i suoi
  operandi esistono, qualunque sia la sua posizione nel programma.
\item \textbf{Conferma, in ordine.} I risultati diventano ufficiali nell'ordine
  del programma: visto da fuori, niente ha mai sorpassato.
\end{enumerate}

L'algoritmo è quello di Tomasulo, progettato nel 1967 per l'unità in virgola
mobile dell'IBM System/360 Model~91; il reorder buffer che conferma in ordine è
arrivato negli anni Ottanta. Il modello del programmatore, quello imparato in
questo libro, è conservato esattamente: è questo il contratto dell'insieme di
istruzioni.

\subsection{La ridenominazione dei registri: trentadue nomi, molti registri}

\asmfile[caption={Lo stesso codice come è scritto, e come è eseguito dopo la ridenominazione.},label=lst:by:rename]{code/l24/c05.S}

Nel Listato~\ref{lst:by:rename} la seconda metà riusa il \emph{nome}
\reg{t0}, non il suo valore: una \term{falsa dipendenza}. Una volta
ridenominate, le due metà non condividono niente, e possono essere eseguite
nello stesso momento su registri diversi. (I numeri da \code{p41} a \code{p44}
sono inventati.) I trentadue registri noti al programmatore sono nomi; il chip
ne ha molti di più, alcune centinaia in un core grande.

\subsection{Indovinare il futuro: la predizione dei salti}

Un processore largo incontra un salto ogni poche istruzioni e non può aspettare
di conoscerne l'esito. Ne \emph{predice} la direzione dalla storia di quel
salto, e continua a eseguire lungo l'ipotesi: un'ipotesi giusta non costa
niente; una sbagliata butta via il lavoro speculativo e ricomincia. Il ciclo
del riscaldamento salta indietro 959 volte e prosegue una volta: un predittore
lo impara nei primi giri. I predittori moderni hanno ragione molto più di nove
volte su dieci.

\begin{prgaside}[Speculazione e sicurezza]
L'esecuzione speculativa ha un lato oscuro. Nel 2018 gli attacchi chiamati
Spectre e Meltdown hanno mostrato che il lavoro buttato via può lasciare
comunque nella cache tracce che rivelano segreti. I trucchi per le prestazioni
sono anche questioni di sicurezza.
\end{prgaside}

\subsection{RISC-V diventa superscalare}

\begin{table}[htb]
\centering\small
\caption{Alcuni core \riscv. L'elenco invecchia in fretta.}\label{tab:by:cores}
\begin{tabularx}{\textwidth}{@{}lP{0.3\textwidth}L@{}}
\toprule
\textbf{Core} & \textbf{Da} & \textbf{Tipo}\\
\midrule
core HP dell'ESP32-C6 & Espressif & in ordine, quattro stadi, un'istruzione per ciclo: il kit\\
BOOM & University of California, Berkeley & fuori ordine, open source, per ricerca e didattica\\
XiangShan & Accademia Cinese delle Scienze & fuori ordine, open source, alte prestazioni\\
XuanTie C910 & T-Head (Alibaba) & fuori ordine, largo tre; rilasciato come open source\\
Performance P550 & SiFive & fuori ordine, largo tre, tredici stadi; venduto come progetto\\
Veyron & Ventana Micro Systems & fuori ordine, per i centri di calcolo\\
\bottomrule
\end{tabularx}
\end{table}

Poiché l'insieme di istruzioni è aperto, chiunque può costruire un core di
questo tipo (Tabella~\ref{tab:by:cores}), e diversi sono pubblicati come codice
sorgente che uno studente può leggere: non ha precedenti tra le architetture
commerciali.

\begin{definition}{Architettura e microarchitettura}
L'\term{architettura} è il contratto: le istruzioni, i registri, che cosa fa
ciascuna istruzione. La \term{microarchitettura} è un modo di onorarlo: stadi,
larghezza, predittori, cache.
\end{definition}

La \ins{add} scritta per un microcontrollore a 160~MHz è la \ins{add} di un
core per server, e \vermone girerebbe senza modifiche su uno qualsiasi dei core
della tabella, dato uno schermo. È per questo che un insieme di istruzioni
imparato una volta è una conoscenza che dura. La distinzione organizza l'intero
libro: i Capitoli da \ref{ch:firstcontact} a~\ref{ch:verm1p4} hanno insegnato
l'architettura, e il kit è una piccola microarchitettura.

% ----------------------------------------------------------------------------
\section{Vettori: un'istruzione, molti dati}

\subsection{La tassonomia di Flynn}

\begin{table}[htb]
\centering\small
\caption{La tassonomia di \citet{flynn1972}.}\label{tab:by:flynn}
\begin{tabular}{@{}llll@{}}
\toprule
\textbf{Classe} & \textbf{Istruzioni} & \textbf{Dati} & \textbf{Esempio}\\
\midrule
SISD & un flusso & un flusso & il kit: un core, un dato alla volta\\
SIMD & un flusso & molti insieme & istruzioni vettoriali; una GPU, in prima approssimazione\\
MISD & molti & un flusso & rara: sistemi ridondanti che votano\\
MIMD & molti & molti & processori multicore, cluster, supercalcolatori\\
\bottomrule
\end{tabular}
\end{table}

La classificazione della Tabella~\ref{tab:by:flynn} è grossolana, ed è ancora
il primo vocabolario del calcolo parallelo.

\subsection{Lo stesso ciclo, con i vettori}

\asmfile[caption={Svuotare la mappa con l'estensione vettoriale.},label=lst:by:vec]{code/l24/c06.S}

L'estensione V di \riscv, ratificata nel 2021, aggiunge i \term{registri
vettoriali}, larghi molti byte (Listato~\ref{lst:by:vec}). L'istruzione
\ins{vsetvli} chiede: quanti dei byte che restano stanno in un registro? La
risposta torna in \reg{t2}. Poi \ins{vse8.v} memorizza \reg{t2} byte con una
sola istruzione. Con registri di 128 bit il ciclo tratta 16 byte per giro: 60
giri invece di 960.

Il codice non dice mai quanto sono larghi i registri: è \emph{agnostico
rispetto alla lunghezza del vettore}, e lo stesso binario usa 128, 256 o
1\,024 bit. L'estensione non è sul kit, dato che l'ESP32-C6 è RV32IMAC; il
listato è stato provato su QEMU con registri di 128, 256, 512 e 1\,024 bit.

\subsection{Una somma, corsia per corsia}

\asmfile[caption={La somma di sedici parole con i vettori.},label=lst:by:vsum]{code/l24/c07.S}

Il Listato~\ref{lst:by:vsum} è la somma parallela dei Capitoli
\ref{ch:information} e~\ref{ch:video}, in \riscv. \ins{vle32.v} carica più
parole insieme; \ins{vadd.vv} somma registro a registro, \emph{corsia} per
corsia (\emph{lane}); alla fine \ins{vredsum.vs} somma tra loro le corsie. Con
i sedici numeri da 1 a 16 il risultato è 136. (L'impostazione \code{tu},
\emph{tail undisturbed}, conserva le somme parziali delle corsie che l'ultimo
giro non usa.)

% ----------------------------------------------------------------------------
\section{Calcolo ad alte prestazioni}

\subsection{Che cos'è un supercalcolatore?}

Un supercalcolatore non è un solo enorme processore. Sono migliaia di
calcolatori ordinari, i \term{nodi}, che lavorano allo stesso problema
(Figura~\ref{fig:by:cluster}). Ogni nodo ha uno o due processori con decine di
core, spesso più GPU, e moltissima memoria. Una rete molto veloce unisce i
nodi, e un file system parallelo li serve tutti. La velocità si misura in
operazioni in virgola mobile al secondo, FLOPS, su un problema standard: un
enorme sistema di equazioni lineari (il benchmark chiamato HPL,
\emph{High-Performance Linpack}). La lista TOP500 classifica le macchine più
veloci due volte l'anno dal 1993.

\begin{figure}[htb]
\centering
\includegraphics[width=0.55\textwidth]{cluster_1.jpg}
\caption{Un piccolo cluster: nodi, GPU, una rete veloce. La struttura di un
supercalcolatore, in piccolo.}
\label{fig:by:cluster}
\end{figure}

\begin{table}[htb]
\centering\small
\caption{Ordini di grandezza (cifre HPL arrotondate dalle liste TOP500; la
cifra per il portatile è un ordine di grandezza).}\label{tab:by:flops}
\begin{tabularx}{\textwidth}{@{}llL@{}}
\toprule
\textbf{Macchina} & \textbf{Dove} & \textbf{Operazioni in virgola mobile al secondo}\\
\midrule
ESP32-C6 & il kit & nessuna unità in virgola mobile: virgola fissa, come in \vermone\\
Un portatile & la vostra borsa & circa $10^{11}$\\
Leonardo (2022) & CINECA, Bologna & circa $2,4\times 10^{17}$\\
Frontier (2022) & Oak Ridge, USA & oltre $10^{18}$: la prima macchina exascale\\
El Capitan (2024) & Livermore, USA & circa $1,7\times 10^{18}$\\
JUPITER (2025) & J\"ulich, Germania & circa $10^{18}$: la prima macchina exascale europea\\
\bottomrule
\end{tabularx}
\end{table}

\emph{Exa} è $10^{18}$. Sette ordini di grandezza separano un portatile da una
macchina exascale (Tabella~\ref{tab:by:flops}): circa cinquant'anni di legge
di Moore, comprati invece con il parallelismo.

\subsection{La legge di Amdahl}

\begin{theorem}{Legge di Amdahl}
Se una frazione $p$ del lavoro di un programma può essere svolta in parallelo,
lo speed-up ottenuto con $n$ processori è
\[ S(n)=\frac{1}{(1-p)+\dfrac{p}{n}},\qquad\text{e}\qquad
   \lim_{n\to\infty}S(n)=\frac{1}{1-p}. \]
\end{theorem}

\noindent\emph{Dimostrazione.} Si prenda come 1 il tempo su un processore. La
parte seriale impiega $1-p$ qualunque sia il numero di processori; la parte
parallela, divisa tra $n$, impiega $p/n$. Lo speed-up è il rapporto tra i due
tempi. \hfill$\square$

\smallskip
La frazione seriale non si riduce, per quanti processori si aggiungano
\citep{amdahl1967}. Con $p=0,9$ lo speed-up non può mai superare 10, nemmeno
con un milione di processori (Tabella~\ref{tab:by:amdahl}): mille processori
su un programma parallelo al 90\% lo rendono meno di dieci volte più veloce.
Prima di comprare processori, misurate $p$. La stessa legge vale per qualsiasi
miglioramento che si applichi solo a una parte del lavoro: una GPU più veloce
non aiuta la parte che gira sulla CPU.

\begin{table}[htb]
\centering\small
\caption{La legge di Amdahl in numeri: lo speed-up.}\label{tab:by:amdahl}
\begin{tabular}{@{}ccccc@{}}
\toprule
\textbf{Frazione parallela $p$} & $n=10$ & $n=100$ & $n=1000$ & $n$ illimitato\\
\midrule
0,50 & 1,8 & 2,0 & 2,0 & 2\\
0,90 & 5,3 & 9,2 & 9,9 & 10\\
0,99 & 9,2 & 50 & 91 & 100\\
0,999 & 9,9 & 91 & 500 & 1000\\
\bottomrule
\end{tabular}
\end{table}

Come verifica, $p=0,9$ e $n=100$ danno $1/(0,1+0,009)=9,17$. La risposta di
Gustafson, nel 1988, fu che con più processori si risolvono problemi più
grandi, e che in un problema più grande la parte parallela di solito cresce
più in fretta di quella seriale.

\subsection{Due modi di dividere il lavoro}

Con la \term{memoria condivisa}, più core di un nodo vedono la stessa memoria,
e il programma gira come \emph{thread} che leggono e scrivono le stesse
variabili. È facile cominciare: OpenMP aggiunge una riga sopra un ciclo. Il
pericolo sono due thread che scrivono la stessa variabile. Con la
\term{memoria distribuita} ogni nodo ha la propria memoria e niente è
condiviso; il programma gira come \emph{processi} che si scambiano messaggi, e
MPI, la \emph{Message Passing Interface}, è lo standard. Scala fino a un
milione di core, e il programmatore ne paga il prezzo. Le applicazioni reali
usano entrambi: MPI tra i nodi, thread e GPU dentro ogni nodo. I modelli
meteorologici e marini del Dipartimento funzionano in questo modo.

\subsection{Il limite è l'energia}

Una macchina exascale assorbe circa venti megawatt, il consumo di una piccola
città (Frontier ha raggiunto il suo exaflop a circa 21~MW): circa 50 miliardi
di operazioni per joule. L'elettricità della sua vita costa quanto la macchina.
Quindi la cifra che conta sono le operazioni per watt, che ha la sua
classifica, la Green500, guidata dalle GPU e da altri acceleratori: l'hardware
specializzato fa di più con ogni joule. La stessa ragione mette nel kit un
piccolo core \riscv: milliwatt. L'efficienza energetica è il filo che lega un
microcontrollore a un supercalcolatore.

\subsection{RISC-V e il supercalcolo europeo}

Le macchine più grandi di oggi usano processori x86 o Arm con acceleratori
GPU. L'Europa progetta i propri processori e acceleratori su \riscv: la
European Processor Initiative, e DARE, il progetto nominato nel
Capitolo~\ref{ch:intro} \citep{eurohpcdare}. Monte Cimone, costruito nel 2022
all'Università di Bologna con CINECA, è un cluster HPC completo su processori
\riscv; il Barcelona Supercomputing Center progetta e fabbrica core \riscv
propri. Un insieme di istruzioni aperto permette a un continente di costruire
ciò che altrimenti dovrebbe comprare.

% ----------------------------------------------------------------------------
\section{Acceleratori GPU}

\subsection{Due modi di spendere i transistor}

Una CPU ha pochi core grandi, alcune decine, ciascuno fuori ordine, con
predittori e grandi cache: è costruita per finire \emph{un} compito il prima
possibile, ed è a suo agio con i salti e con i dati irregolari. Una GPU ha
migliaia di core piccoli e semplici, che eseguono tutti lo stesso codice su
dati diversi: è costruita per finire un milione di compiti al secondo, ed è a
suo agio con gli array regolari: pixel, matrici, griglie. Sono i buoi e i polli
di Cray, e la risposta dipende dal campo. La GPU non è più veloce in nessuna
cosa singola; fa più cose insieme.

\subsection{Come un programma usa una GPU}

\begin{enumerate}
\item \textbf{Preparare.} Il programma sulla CPU, l'\emph{host}, costruisce i
  suoi array in memoria centrale.
\item \textbf{Copiare dentro.} I dati viaggiano verso la memoria della GPU: un
  costo reale, in tempo.
\item \textbf{Lanciare.} Una piccola funzione, il \term{kernel}, gira in
  migliaia di thread, uno per elemento dei dati.
\item \textbf{Copiare fuori.} I risultati tornano indietro, e l'host prosegue.
\end{enumerate}

La GPU è un dispositivo, come quelli del Capitolo~\ref{ch:io}: la CPU le invia
dati e comandi attraverso un driver. I modelli di programmazione sono CUDA,
OpenCL e SYCL, e direttive come OpenACC e l'offload di OpenMP.

\cfile[caption={La mappa di \vermone con un thread per cella (CUDA C, nella
sua forma usuale da libro di testo; non compilato per questo libro).},label=lst:by:cuda]{code/l24/c08.c}

Nel Listato~\ref{lst:by:cuda} non c'è alcun ciclo: ci sono 960 thread. Ogni
thread calcola il proprio indice e fa uno store. I thread arrivano a blocchi,
e gli ultimi non trovano niente da fare. Per 960 byte le due copie costano
molto più del lavoro, e il punto è serio: una GPU conviene quando i dati sono
grandi e il lavoro per elemento è reale.

\subsection{Quando conviene una GPU?}

Conviene quando la stessa operazione si applica a moltissimi dati, in modo
indipendente: immagini, matrici, le griglie di una simulazione, reti neurali; e
quando i dati possono restare sulla GPU per molti passi invece di viaggiare a
ognuno. Non conviene quando il lavoro è una catena di decisioni, ciascuna
dipendente dalla precedente. E vale di nuovo la legge di Amdahl: ciò che resta
sulla CPU fissa il limite. Una simulazione meteorologica aggiorna ogni cella di
una griglia tridimensionale a partire dalle sue vicine, milioni di volte: la
stessa forma della mappa di \vermone, un po' più grande. L'addestramento delle
reti neurali di oggi è, quasi interamente, aritmetica matriciale su GPU.

% ----------------------------------------------------------------------------
\section{Acceleratori FPGA}

\subsection{Un chip che diventa il circuito che descrivete}

\begin{definition}{FPGA}
Una \term{FPGA} (\emph{field-programmable gate array}) è un chip fatto di un
mare di piccoli blocchi programmabili: \term{tabelle di look-up}, ciascuna
capace di essere una qualsiasi funzione di pochi ingressi; flip-flop, uno
accanto a ogni tabella; e fili e interruttori che possono unire qualsiasi
blocco a qualsiasi altro. Un file di configurazione, il \term{bitstream},
imposta ogni tabella e ogni interruttore.
\end{definition}

Caricate un altro bitstream, e lo stesso chip è un altro circuito. La prima
FPGA, la Xilinx XC2064, fu venduta nel 1985; oggi vanno da pochi euro a schede
per centri di calcolo. Una FPGA non esegue un programma: \emph{diventa} un
circuito. Niente viene prelevato o decodificato, e tutte le porte lavorano
nello stesso momento.

\subsection{Una tabella di look-up è una tabella di verità in memoria}

\begin{center}\small\ttfamily
\begin{tabular}{cc|ccc}
\textrm{$a$}&\textrm{$b$}&\textrm{memorizzato per AND}&\textrm{memorizzato per XOR}&\textrm{memorizzato per NOT $a$}\\\hline
0&0&0&0&1\\0&1&0&1&1\\1&0&0&1&0\\1&1&1&0&0
\end{tabular}
\end{center}

Gli ingressi sono l'indirizzo, e quattro bit memorizzati scelgono la funzione:
cambiare i bit cambia la porta. Il Capitolo~\ref{ch:boole1} diceva che una
tabella di verità definisce una funzione, e il Capitolo~\ref{ch:storage} che
una ROM indirizzata dagli ingressi realizza qualsiasi funzione. Una tabella di
look-up è quella ROM, con quattro-sei ingressi e da sedici a sessantaquattro
bit; il flip-flop accanto a essa è quello del Capitolo~\ref{ch:sequential}.
Una FPGA è questo libro, ripetuto centomila volte su un chip.

\subsection{La regola dell'inversione a U di Verm1, come hardware}

\lstinputlisting[style=verilog,caption={La regola dell'inversione a U in Verilog (non sintetizzata per questo libro).},label=lst:by:verilog]{code/l24/c09.v}

Verilog, come VHDL, è un \term{linguaggio di descrizione dell'hardware}:
descrive circuiti, non passi (Listato~\ref{lst:by:verilog}). L'istruzione
\code{assign} è un filo, sempre vera: non c'è un prima e non c'è un dopo. Il
circuito è fatto di due porte XOR e un AND con un ingresso invertito,
\[ \mathit{refuse}=(d_1\oplus w_1)\cdot\nott{(d_0\oplus w_0)} . \]
Sul kit questa regola richiedeva tre istruzioni e tre cicli
(Capitolo~\ref{ch:verm1p2}); qui le porte rispondono in una frazione di
nanosecondo, in continuazione.

\subsection{Dalla descrizione al circuito}

\begin{figure}[htb]
\centering
\begin{tikzpicture}[node distance=4mm]
  \node[blkacc,text width=21mm,minimum height=16mm] (a) {\textbf{Descrivere}\\[1pt]\footnotesize Verilog, VHDL, o C};
  \node[blkamb,text width=21mm,minimum height=16mm,right=of a] (b) {\textbf{Simulare}\\[1pt]\footnotesize forme d'onda};
  \node[blkvio,text width=21mm,minimum height=16mm,right=of b] (c) {\textbf{Sintetizzare}\\[1pt]\footnotesize tabelle e flip-flop};
  \node[blkgrn,text width=21mm,minimum height=16mm,right=of c] (d) {\textbf{Piazzare e instradare}\\[1pt]\footnotesize sul chip};
  \node[blkred,text width=21mm,minimum height=16mm,right=of d] (e) {\textbf{Configurare}\\[1pt]\footnotesize caricare il bitstream};
  \foreach \x/\y in {a/b,b/c,c/d,d/e} \draw[->,thick] (\x) -- (\y);
\end{tikzpicture}
\caption{Da una descrizione a un circuito.}
\label{fig:by:fpgaflow}
\end{figure}

Ogni tentativo richiede da minuti a ore (Figura~\ref{fig:by:fpgaflow}): più
lento che compilare una cartuccia, molto più veloce che fabbricare un chip. Lo
stesso flusso, con una fabbrica alla fine invece di un bitstream, produce un
\term{ASIC}, un chip fatto per un solo scopo. Le FPGA sono il luogo in cui i
chip si provano prima di fabbricarli.

\subsection{Anche un processore è un circuito}

Un core \riscv è una descrizione in un linguaggio hardware, e può essere
caricato su una FPGA. Questi \term{soft core} sono liberi da leggere e da
usare: PicoRV32, VexRiscv, NEORV32, Rocket. Un core della classe di quello del
kit sta in una FPGA piccola ed economica, e accanto a esso si può mettere un
circuito proprio per l'unica operazione che un programma ripete di più: un
acceleratore su misura. \riscv riserva spazio nella sua codifica per istruzioni
personalizzate, con cui raggiungere tali circuiti.

\begin{prgconcept}[Il cerchio si chiude]
Le porte del Capitolo~\ref{ch:boole1} costruiscono il sommatore del
Capitolo~\ref{ch:boole2}, i registri del Capitolo~\ref{ch:registers} e il
processore del Capitolo~\ref{ch:processor}; e quel processore è un file che si
può leggere. Dalle porte a un processore, e ritorno: ora conoscete ogni strato
intermedio.
\end{prgconcept}

\subsection{Quattro modi di calcolare}

\begin{table}[htb]
\centering\small
\caption{Quattro modi di calcolare.}\label{tab:by:four}
\begin{tabularx}{\textwidth}{@{}lLLLL@{}}
\toprule
 & \textbf{CPU} & \textbf{GPU} & \textbf{FPGA} & \textbf{ASIC}\\
\midrule
Che cos'è & pochi core complessi & migliaia di core semplici & un circuito che si configura & un circuito fatto per un solo lavoro\\
Si programma in & qualsiasi linguaggio & kernel: CUDA, OpenCL, SYCL & Verilog, VHDL, HLS & gli stessi, poi una fabbrica\\
Forte in & tutto, moderatamente & dati regolari in massa & flussi, bassa latenza & una cosa, alla perfezione\\
Lavoro per joule & basso & alto & alto & il più alto\\
Sforzo d'uso & basso & medio & alto & molto alto\\
Cambiare idea & ricompilare & ricompilare & un nuovo bitstream & un nuovo chip\\
\bottomrule
\end{tabularx}
\end{table}

La generalità si paga in energia; l'efficienza si paga in sforzo
(Tabella~\ref{tab:by:four}). Un moderno sistema su un chip li ha tutti e
quattro: un telefono contiene core CPU, una GPU, e circuiti fissi per il video
e la radio.

% ----------------------------------------------------------------------------
\section{Ritorno alla macchina}

\begin{table}[htb]
\centering\small
\caption{Svuotare la mappa di \vermone, 960 celle, su cinque macchine.}\label{tab:by:five}
\begin{tabularx}{\textwidth}{@{}P{0.27\textwidth}LP{0.22\textwidth}@{}}
\toprule
\textbf{Macchina} & \textbf{Come} & \textbf{Passi nel tempo}\\
\midrule
Il kit: ESP32-C6 & un byte per giro, un'istruzione per ciclo & 3\,840 cicli\\
Core \riscv fuori ordine & lo stesso binario; più istruzioni per ciclo, il salto predetto & una frazione di quei cicli\\
\riscv con vettori, 128 bit & 16 byte per giro & 60 giri\\
GPU & 960 thread, uno per cella & un lancio, più due copie\\
FPGA & 960 flip-flop con un reset comune & un ciclo di clock\\
\bottomrule
\end{tabularx}
\end{table}

Le prime tre righe della Tabella~\ref{tab:by:five} sono lo stesso insieme di
istruzioni; la prima e la terza sono state eseguite per questo libro. Il
lavoro è lo stesso dappertutto; ciò che cambia è quanto di esso avviene nello
stesso momento. E cambiano i costi: alla prima riga servono un chip da due
euro e quattro righe di assembly; all'ultima serve un progetto hardware.

\begin{prgconcept}
La domanda dell'ingegnere non è mai quale macchina sia la più veloce, ma quale
basti.
\end{prgconcept}

% ----------------------------------------------------------------------------
% Generated by tools/textbook/export_listings.js from decks/plan.js and code/l24.S: do not edit.
\section{Il programma di questo capitolo}\label{sec:lab:l24}

I listati di questo capitolo sono parti di un solo programma,
\fname{code/l24.S} nel repository del corso, stampato per intero nel
Listato~\ref{lst:lab:l24}. I test automatici del materiale del corso lo
compilano e lo eseguono sull'esecutore senza interfaccia di \prgqt{}, con
ingressi programmati, e controllano il frame; la
Figura~\ref{fig:lab:l24} è il frame scritto dal test.

\begin{description}[style=nextline,leftmargin=1.5em]
\item[Sullo schermo.] Un quadrato verde: i 960 byte della mappa, riempiti di uni da init, sono stati tutti posti a zero dal ciclo di quattro istruzioni della lezione. Un quadrato rosso vorrebbe dire che qualche byte è sopravvissuto.
\item[Che cosa premere.] Niente: il programma va da solo.
\end{description}

\begin{figure}[htb]
\centering
\includegraphics[width=0.5\textwidth]{lab_l24.png}
\caption{\fname{l24.S} su \prgqt{}.}\label{fig:lab:l24}
\end{figure}

\asmfile[caption={Il programma completo del capitolo, \fname{code/l24.S}.},label=lst:lab:l24]{code/l24/program.S}

Per eseguirlo su \prgqt{} i comandi sono quelli del
Listato~\ref{lst:lab:l24:run}; il Paragrafo~\ref{sec:prg:testing} dà gli
stessi passi per QEMU e per la scheda.

\begin{shell}[caption={Compilare \fname{l24.S} ed eseguirlo su \prgqt{}.},label=lst:lab:l24:run]
# 1. start PRG32-QT on this computer (it listens on port 8080)

# 2. in the PRG32 folder, with ESP-IDF active: build the cartridge
python3 -m prg32 cartridge build ~/arc12/code/l24.S --portable \
    --entry-prefix l24 --name l24 --out l24.prg32

# 3. send it to PRG32-QT and start it
python3 -m prg32 esp32c6 upload-and-run l24.prg32 \
    --url http://localhost:8080 --slot cart0

# 4. look, press, compare with the previous slide
\end{shell}


% ----------------------------------------------------------------------------
\section{Verificate di aver capito}

\begin{quiz}{Un programma è parallelo all'80\%. Qual è il massimo speed-up che un numero qualunque di processori può dare?}
\item 4
\item 5
\item 8
\item 80
\end{quiz}

\begin{quiz}{Che cosa elimina la ridenominazione dei registri?}
\item i salti
\item il bisogno di uno stack
\item le false dipendenze tra usi dello stesso nome di registro
\item i cache miss
\end{quiz}

\begin{quiz}{Con registri vettoriali a 256 bit, quanti giri servono al ciclo vettoriale per i 960 byte?}
\item 4
\item 30
\item 60
\item 960
\end{quiz}

\begin{quiz}{Una tabella di look-up a due ingressi memorizza, per gli ingressi 00, 01, 10, 11, i bit 0, 1, 1, 1. Quale porta è?}
\item AND
\item OR
\item XOR
\item NAND
\end{quiz}

\section{Esercizi}

\begin{exercise}[prima su carta, poi sulla macchina]
\begin{enumerate}[nosep]
\item Compilate la tabella di Amdahl per $p=0,95$ e $n=10$, 100, 1000.
\item Disegnate il diagramma della pipeline di un processore largo tre per nove
  istruzioni indipendenti con cinque stadi.
\item Ridenominate i registri del codice di \code{step} che sposta e fa
  rientrare la testa in \vermone, parte~2: quali istruzioni potrebbero essere
  eseguite nello stesso ciclo?
\item Scrivete la tabella di verità memorizzata in una tabella di look-up a
  quattro ingressi per
  $\mathit{refuse}=(\mathit{dir}\oplus\mathit{want})=2$.
\item Compilate il test \fname{vector\_qemu.S} del materiale del corso ed
  eseguitelo su QEMU con registri vettoriali di 128 e di 512 bit: perché lo
  stesso binario funziona?
\end{enumerate}
\expected{Amdahl con $p=0,95$: 6,9, 16,8, 19,6. Nove istruzioni su un
processore largo tre: sette cicli. La tabella di look-up ha quattro uni su
sedici elementi, per $(\mathit{dir},\mathit{want})=(0,2)$, $(1,3)$, $(2,0)$,
$(3,1)$. QEMU si avvia con
\code{qemu-system-riscv32 -machine virt -cpu rv32,v=true,vlen=128}.}
\end{exercise}

\section{Glossario}

\begin{glossarytable}
superscalar & superscalare\\
register renaming & ridenominazione dei registri\\
out-of-order execution & esecuzione fuori ordine\\
branch prediction & predizione dei salti\\
speed-up & accelerazione (speed-up)\\
node & nodo\\
shared memory & memoria condivisa\\
message passing & scambio di messaggi\\
accelerator & acceleratore\\
kernel (GPU) & kernel\\
look-up table & tabella di look-up\\
bitstream & bitstream\\
hardware description language & linguaggio di descrizione dell'hardware\\
soft core & soft core\\
\end{glossarytable}

\section{Conclusioni e letture}

Un'equazione, tre fattori: ogni macchina veloce ne attacca uno, o moltiplica i
processori. I core superscalari fuori ordine rispettano il contratto
dell'insieme di istruzioni ed eseguono molte istruzioni insieme. Vettori, GPU e
cluster sfruttano i dati che possono essere elaborati in modo indipendente, e
la legge di Amdahl dice fin dove si arriva. Le FPGA trasformano una descrizione
in un circuito: la logica di questo libro, in grande scala.

Il libro finisce qui: da un bit a un supercalcolatore, con un solo insieme di
istruzioni.

\paragraph{Letture.} \citet{patterson2020}, capitoli 4 e~6;
\citet{hennessy2019}; \url{https://top500.org}; la specifica vettoriale di
\riscv su \url{https://riscv.org}.
